\section{Environment Details}
\label{app:env}

\paragraph{Transitions.}
Given the current true value $v$, a transition is drawn uniformly from the available operations: \texttt{add}~$k$ with $k \in \{1,\dots,4\}$; \texttt{sub}~$k$ with $k \in \{1,\dots,\min(4, v-1)\}$ if $v > 2$; and \texttt{mul}~$k$ with $k \in \{2,3\}$ if $|v| \le 16$. The corrupted version of \texttt{add}/\texttt{sub}~$k$ uses $k{+}1$, and that of \texttt{mul}~$k$ is \texttt{add}~$k$. If the corrupted operation happens to give the same value, it is replaced by \texttt{add}~$k{+}2$. A corrupted step therefore always changes the value (unit-tested).

\paragraph{Fault prior.}
Conditional on a fault, its position $h \in \{0,\dots,d-1\}$ has weight $1 + 5\,(1 - |h - c| / \max(1, c))$ with $c = (d-1)/2$, i.e.\ it peaks in the middle of the chain.

\paragraph{Cue selectivity metric.}
At every decision point where the current step $t > 0$ has not been checked and \texttt{VERIFY} is available, we record the cue $c_t$ and whether the agent chose \texttt{VERIFY}. $P(\texttt{VERIFY} \mid c)$ pools these counts over all evaluation episodes of a seed.

\paragraph{Validation against the original implementation.}
The environment was rewritten for this version. Over 86{,}112 random steps from 18{,}000 episodes, covering all reward schemes and both budget-penalty settings, the rewritten environment reproduces the original implementation's transitions, masks, rewards and absolute observations exactly. The only difference is the observation returned after the terminal \texttt{COMMIT}, which no learner uses.

\section{Hyperparameters and Compute}
\label{app:hparams}

\begin{table}[h]
\centering
\small
\setlength{\tabcolsep}{4pt}
\caption{Default environment and learner settings.}
\label{tab:hparams}
\begin{tabular}{ll}
\toprule
Setting & Value \\
\midrule
Training / held-out lengths & 3--6 / 7--9 and 10--14 \\
Fault probability; cue hit / false-positive rate & 0.8; 0.85 / 0.15 \\
Budget; action costs (\texttt{THINK}, \texttt{VERIFY}, \texttt{REVISE}) & $d + 6$; $(1,1,1)$ \\
Rewards: success / failure / per-unit cost & $+1$ / $-1.2$ / $-0.04$ \\
Monitorability-aware coefficients (Eq.~\ref{eq:ma}) & $0.02, 0.25, 0.18, 0.02, 0.06, 0.30$ \\
Network & MLP, 2 hidden layers of 64 ReLU units \\
Optimiser; learning rate; gradient clip & Adam; $10^{-3}$; 1.0 (global norm) \\
Discount; loss & 0.99; Huber \\
Replay size; batch size; warm-up & 25{,}000; 64; 500 steps \\
Target network & hard copy every 200 environment steps \\
Exploration & $\epsilon$: 0.5 $\to$ 0.02 linearly over 85\% of training \\
Training episodes per run & 20{,}000 \\
Seeds & 10 (main, ablations), 5 (sensitivity) \\
Validation for selection & 200 episodes of lengths 3--6 every 1{,}000 episodes \\
Evaluation & 500 episodes per seed and setting \\
\bottomrule
\end{tabular}
\end{table}

A training run takes about four minutes on one CPU core. The main experiment (60 runs), the ablations (80 additional runs), the sensitivity study (30 runs) and the follow-up encoder (20 runs) together take about four hours on a 4-core machine. No GPU is used.

\section{Last-Checkpoint Results}
\label{app:final}

\begin{table}[h]
\centering
\small
\setlength{\tabcolsep}{3.5pt}
\caption{As Table~\ref{tab:main}, but evaluating the last checkpoint of each run instead of the one selected on training-length validation. Heuristics are repeated for reference.}
\label{tab:final}
\resizebox{\linewidth}{!}{\begin{tabular}{lccccc}
\toprule
Method & Success 3--6 & Success 7--9 & Success 10--14 & Detect 7--9 & Cost-adj. 7--9 \\
\midrule
Random & 0.00\,{\scriptsize[0.00, 0.01]} & 0.00\,{\scriptsize[0.00, 0.00]} & 0.00\,{\scriptsize[0.00, 0.00]} & 0.02\,{\scriptsize[0.01, 0.02]} & -0.09\,{\scriptsize[-0.09, -0.08]} \\
Always THINK & 0.20\,{\scriptsize[0.19, 0.21]} & 0.20\,{\scriptsize[0.19, 0.21]} & 0.20\,{\scriptsize[0.19, 0.21]} & 0.00\,{\scriptsize[0.00, 0.00]} & -0.20\,{\scriptsize[-0.21, -0.19]} \\
Verify every step & 0.80\,{\scriptsize[0.79, 0.81]} & 0.06\,{\scriptsize[0.06, 0.07]} & 0.00\,{\scriptsize[0.00, 0.00]} & 0.93\,{\scriptsize[0.92, 0.94]} & -0.63\,{\scriptsize[-0.64, -0.63]} \\
Cue-triggered & 0.87\,{\scriptsize[0.87, 0.88]} & 0.88\,{\scriptsize[0.87, 0.89]} & 0.85\,{\scriptsize[0.84, 0.86]} & 0.85\,{\scriptsize[0.84, 0.87]} & 0.33\,{\scriptsize[0.31, 0.34]} \\
Budget-aware & 0.97\,{\scriptsize[0.96, 0.97]} & 0.82\,{\scriptsize[0.81, 0.84]} & 0.68\,{\scriptsize[0.66, 0.70]} & 0.78\,{\scriptsize[0.76, 0.79]} & 0.19\,{\scriptsize[0.17, 0.20]} \\
\midrule
Outcome (inv) & 0.40\,{\scriptsize[0.30, 0.50]} & 0.33\,{\scriptsize[0.24, 0.45]} & 0.30\,{\scriptsize[0.23, 0.40]} & 0.25\,{\scriptsize[0.09, 0.42]} & -0.31\,{\scriptsize[-0.40, -0.19]} \\
Naive process (inv) & 0.41\,{\scriptsize[0.28, 0.56]} & 0.36\,{\scriptsize[0.24, 0.50]} & 0.31\,{\scriptsize[0.23, 0.41]} & 0.58\,{\scriptsize[0.44, 0.71]} & -0.30\,{\scriptsize[-0.42, -0.16]} \\
Monitorability (inv) & 0.77\,{\scriptsize[0.70, 0.84]} & 0.63\,{\scriptsize[0.53, 0.73]} & 0.52\,{\scriptsize[0.43, 0.63]} & 0.81\,{\scriptsize[0.75, 0.86]} & -0.02\,{\scriptsize[-0.13, 0.07]} \\
Outcome (abs) & 0.24\,{\scriptsize[0.22, 0.25]} & 0.12\,{\scriptsize[0.06, 0.18]} & 0.08\,{\scriptsize[0.03, 0.13]} & 0.06\,{\scriptsize[0.01, 0.16]} & -0.20\,{\scriptsize[-0.32, -0.09]} \\
Naive process (abs) & 0.23\,{\scriptsize[0.17, 0.28]} & 0.12\,{\scriptsize[0.06, 0.18]} & 0.05\,{\scriptsize[0.00, 0.10]} & 0.19\,{\scriptsize[0.07, 0.33]} & -0.28\,{\scriptsize[-0.39, -0.17]} \\
Monitorability (abs) & 0.63\,{\scriptsize[0.59, 0.67]} & 0.24\,{\scriptsize[0.16, 0.31]} & 0.03\,{\scriptsize[0.00, 0.07]} & 0.51\,{\scriptsize[0.38, 0.61]} & -0.34\,{\scriptsize[-0.42, -0.27]} \\
\bottomrule
\end{tabular}
}
\end{table}

\section{Changes from the Earlier Version}
\label{app:changes}

The first public version of this repository reported that monitorability-aware RL reached 0.60 success on training lengths but only 0.11 on lengths 7--9, and concluded that the learned checking behaviour was brittle under length shift. This version changes the following.
\begin{enumerate}
  \item \textbf{Representation control.} The original observation contained absolute position and depth features. We add a length-invariant encoder and compare the two directly (\S\ref{sec:length}). The encoder, not the reward, explains most of the length-generalization failure.
  \item \textbf{Training length and checkpoint selection.} Runs were trained for 5{,}000 episodes, while learning curves were still rising, and evaluated at the last checkpoint. We train for 20{,}000 episodes, select checkpoints on training-length validation only, and report last-checkpoint results as well (Appendix~\ref{app:final}).
  \item \textbf{Evaluation leakage.} Training episode seeds were \texttt{seed}$\cdot 10^5 + i$ and evaluation seeds \texttt{seed}$\cdot 10^6 + i$, which coincide for seed~0: 299 of seed~0's 300 training-length evaluation episodes were exact copies of training episodes. Training, validation and evaluation now use disjoint random-number streams, and a regression test checks that they stay disjoint.
  \item \textbf{Localization metric.} The original ``failed-episode localization score'' gave full credit when the fault was on the last step and the agent committed without checking anything; as a result the never-check policy scored 0.07. We replace it with fault detection, the probability that a check caught the fault.
  \item \textbf{Stronger baselines and statistics.} We add a budget-aware heuristic tuned on training lengths, use 10 seeds instead of 5, and report bootstrap confidence intervals and paired differences instead of standard errors.
  \item \textbf{Heuristic definitions.} The original cue-triggered heuristic re-checked a step after repairing it; ours does not, which raises its success from 0.83 to 0.88. Learned agents are unaffected.
\end{enumerate}
